\documentclass[10pt,a4paper]{amsart}
\usepackage{amsmath,amssymb,amsthm,graphicx}
\newcommand{\newq}{\mathop{\mathsf N}}
\AtBeginDocument{\pagestyle{plain}}
\title{Generalised fresh-quantified permutation fixed points decompose into freshness and algebraic invariance}
\author{Ali K. Caires-Santos, Maribel Fernández, Murdoch James Gabbay, Daniele Nantes-Sobrinho}
\date{doi:10.4230/LIPIcs.FSCD.2026.10; CC BY4.0 (official publisher article)}
\begin{document}
\maketitle
\noindent\textit{Editorial note.} Source https://drops.dagstuhl.de/entities/document/10.4230/LIPIcs.FSCD.2026.10. Copyright Ali K. Caires-Santos, Maribel Fernández, Murdoch James Gabbay and Daniele Nantes-Sobrinho, CC BY4.0. Theorem2.3.5 and complete author appendix proof/support faithfully transcribed. Original shorthand references and notation are retained without added assumptions or proofs.\par\medskip
\noindent\textit{Allowed standard prerequisites (editorial).}\begin{itemize}
\item Finite permutation groups, support equivariance and the single-atom Some/Any principle. Publisher sections2.1-2.3 and cited Proposition4.10 [16]. Actual nested Some/Any and cycle-support proofs are included, as is the whole two-part generalisation.
\end{itemize}\medskip
\section*{Permutations, support and fresh quantification (publisher sections 2.1--2.3)}
We fix disjoint countably infinite sets: the set of atoms $\mathbb A=\{a,b,c,\ldots\}$ and unknowns $\mathbb V=\{X,Y,Z,\ldots\}$.
A (finite) permutation on $\mathbb A$ is a bijection $\pi:\mathbb A\to\mathbb A$ whose set of moved atoms $\operatorname{mv}(\pi)=\{a\in\mathbb A:\pi(a)\ne a\}$ is finite. Write $\pi\circ\pi'$, $\mathsf{id}$ and $\pi^{-1}$ for functional composition, the identity, and the inverse of permutations, respectively.
The group of all finite permutations of atoms is denoted $G(\mathbb A)$.
Every finite permutation $\pi\in G(\mathbb A)$ can be uniquely decomposed (up to reordering) into pairwise disjoint cycles as $\pi=\gamma_1\circ\cdots\circ\gamma_n$, with $\operatorname{mv}(\gamma_i)\cap\operatorname{mv}(\gamma_j)=\varnothing$ for $i\ne j$.
This is called the canonical form of $\pi$.
\paragraph{Assumption 2.1.1.} Permutations are always in a canonical form.
A nominal set is a pair $(\mathcal X,\cdot)$ consisting of a set $\mathcal X$ and an action $\cdot:G(\mathbb A)\times\mathcal X\to\mathcal X$ such that every element $x\in\mathcal X$ has finite support, i.e., there exists a finite $B\subseteq\mathbb A$ with $\pi\cdot x=x$ for all $\pi\in G(\mathbb A)$ such that $\pi(a)=a$ for every $a\in B$.
The least such $B$ is denoted $\operatorname{supp}(x)$. An atom $a$ is fresh for $x$ (written $a\#x$) if $a\notin\operatorname{supp}(x)$.
Sometimes, when there is no ambiguity, we refer to a nominal set $(\mathcal X,\cdot)$ by its carrier set $\mathcal X$ and leave the action $\cdot$ implicit.
Throughout this section, let $\mathcal X$ denote a nominal set. We consider formulas of the form $\phi(c,\bar x)$, where $c$ is an atom and $\bar x=x_1,\ldots,x_n$ is a finite list of parameters with each $x_i\in\mathcal X$.
Typically, $\phi$ is a predicate in higher-order logic [22] or ZFA set theory [14,16].
\paragraph{Definition 2.3.1.} The $\mathsf N$-quantifier is formally defined as:
\[
\newq c.\phi(c,\bar x)\ \Longleftrightarrow\ \{c\in\mathbb A:\phi(c,\bar x)\}\text{ is cofinite}.
\]
A cofinite subset of a set $M$ is a subset $B$ whose complement in $M$ is a finite set.
In cases where multiple fresh atoms are introduced at once, we consider nested $\mathsf N$-quantification of the form $\newq c_1\cdots\newq c_n.\phi(c_1,\ldots,c_n,\bar x)$, abbreviated as $\newq\bar c.\phi(\bar c,\bar x)$.
For a finite set of atoms $A$ and a finite list of elements $\bar x=x_1,\ldots,x_k$ of $\mathcal X$, write $A\#\bar x$ as abbreviation for $A\cap\operatorname{supp}(\bar x)=\varnothing$, where $\operatorname{supp}(\bar x)=\bigcup_{j=1}^k\operatorname{supp}(x_j)$.
If $\bar a=a_1,\ldots,a_n$ is a list of atoms, we write $\bar a\#x$ as shorthand for $\{\bar a\}\#x$, where $\{\bar a\}=\{a_1,\ldots,a_n\}$.
A fundamental connection between the $\mathsf N$-quantifier and freshness is given by the following equivalence [16]:
\[
a\#x\ \Longleftrightarrow\ \newq c.(a\ c)\cdot x=x.
\]
We define a generalised fixed-point identity as an assertion of the form $\newq\bar c.\pi\cdot x=x$, where $\pi$ is a permutation that may contain atoms bound by $\mathsf N$.
\paragraph{Definition 2.3.2 ($\bar c$-cycles and $\bar c$-permutations).}
Let $\pi$ be a permutation written as a product of disjoint cycles and let $\bar c$ be a list of atoms.
\begin{enumerate}
\item A cycle $\eta$ of $\pi$ is called a $\bar c$-cycle if $\operatorname{mv}(\eta)\cap\{\bar c\}\ne\varnothing$, and a $\neg\bar c$-cycle otherwise.
\item A permutation $\pi$ is a $\bar c$-permutation iff all its cycles are $\bar c$-cycles, and a $\neg\bar c$-permutation iff all its cycles are $\neg\bar c$-cycles.
\end{enumerate}
Write $\pi=\pi_{\bar c}\circ\pi_{\neg\bar c}$ where $\pi_{\bar c}$ is the product of all $\bar c$-cycles of $\pi$, and $\pi_{\neg\bar c}$ is the product of all $\neg\bar c$-cycles of $\pi$.
For a set $S\subseteq G(\mathbb A)$, we write $\langle S\rangle$ for the subgroup generated by $S$, that is, the smallest subgroup of $G(\mathbb A)$ containing $S$.
In particular, if $S=\{\rho_1,\ldots,\rho_n\}$ is a finite set of permutations, we write $\langle\rho_1,\ldots,\rho_n\rangle$ for $\langle S\rangle$.

\section*{Complete author supporting proofs (publisher appendix A)}
\paragraph{Lemma 2.3.4.} Let $x\in\mathcal X$ and $\pi,\eta_1,\ldots,\eta_n$ be permutations written as compositions of disjoint cycles.
\begin{enumerate}
\item If $\pi\cdot x=x$, $\eta$ is a cycle of $\pi$, and $\operatorname{mv}(\eta)-\operatorname{supp}(x)\ne\varnothing$, then $\operatorname{mv}(\eta)\#\operatorname{supp}(x)$.
\item If $\eta_i\cdot x=x$, for all $i=1,\ldots,n$, and $\rho\in\langle\eta_1,\ldots,\eta_n\rangle$, then $\rho\cdot x=x$.
\end{enumerate}
\begin{proof}
\begin{enumerate}
\item Write $\pi=\mu_1\circ\cdots\circ\mu_n$ with disjoint cycles $\mu_1,\ldots,\mu_n$, and let $\eta=\mu_j$ for some $1\le j\le n$. Choose $a\in\operatorname{mv}(\eta)\setminus\operatorname{supp}(x)$.
Suppose, by contradiction, there exists $b\in\operatorname{mv}(\eta)\cap\operatorname{supp}(x)$.
Since the cycles are disjoint, $\mu_i(b)=b$ for all $i\ne j$. Hence $\pi(b)=\eta(b)$.
Since $b\in\operatorname{supp}(x)$ and $\pi\cdot x=x$, it follows that $\pi(b)\in\operatorname{supp}(\pi\cdot x)=\operatorname{supp}(x)$, and therefore $\eta(b)\in\operatorname{supp}(x)$.
Because $a,b\in\operatorname{mv}(\eta)$, there exists $l>0$ such that $\eta^l(b)=a$.
Applying $\pi$ repeatedly to $\eta(b)\in\operatorname{supp}(x)$ and using that $\pi\cdot x=x$ yields $\eta^l(b)\in\operatorname{supp}(x)$, hence $a\in\operatorname{supp}(x)$, contradicting the choice of $a$.
\item Assume that $\eta_i\cdot x=x$ for every $i=1,\ldots,n$, and let $\rho\in\langle\eta_1,\ldots,\eta_n\rangle$.
By definition of the subgroup generated by $\eta_1,\ldots,\eta_n$, the permutation $\rho$ can be written as a finite composition $\rho=\zeta_1\circ\cdots\circ\zeta_k$, where each $\zeta_r$ belongs to $\{\eta_1^{\pm1},\ldots,\eta_n^{\pm1}\}$.
We first note that each inverse $\eta_i^{-1}$ also fixes $x$.
Indeed, from $\eta_i\cdot x=x$, applying $\eta_i^{-1}$ to both sides gives $\eta_i^{-1}\cdot(\eta_i\cdot x)=\eta_i^{-1}\cdot x$, and therefore $x=\eta_i^{-1}\cdot x$.
So every $\zeta_r$ in the above decomposition fixes $x$.
We now prove that $\rho$ fixes $x$ by repeated application of these equalities:
\[
\rho\cdot x=(\zeta_1\circ\cdots\circ\zeta_k)\cdot x=\zeta_1\cdot(\zeta_2\cdot(\cdots(\zeta_k\cdot x)\cdots)).
\]
Since $\zeta_k\cdot x=x$, the innermost term simplifies to $x$; then $\zeta_{k-1}\cdot x=x$, and continuing in this way we obtain $\rho\cdot x=x$.
\end{enumerate}
\end{proof}
A key property of the $\mathsf N$-quantifier is the Some/Any principle [16, Prop.\ 4.10]:
\paragraph{Proposition A.1 (Some/Any Principle).} The following are equivalent:
\[
\newq c.\phi(c,\bar x)\ \Longleftrightarrow\ \forall c.(c\#\bar x\Rightarrow\phi(c,\bar x))\ \Longleftrightarrow\ \exists c.(c\#\bar x\land\phi(c,\bar x)).
\]
For multiple atoms, nested $\mathsf N$-quantifiers introduce distinctness conditions. We avoid writing these explicitly by fixing the following convention.
\paragraph{Notation A.2.} Let $\bar c=c_1,\ldots,c_n$ be a list of pairwise distinct atoms.
We abbreviate $\forall c_1\cdots\forall c_n.(\bigwedge_{i<j}c_i\ne c_j\Rightarrow\phi)$ as $\forall\bar c.\phi$, and similarly for $\exists\bar c.\phi$.
\paragraph{Lemma A.3 (Generalised Some/Any Principle).} Let $\bar c$ be a list of atoms. Then:
\[
\newq\bar c.\phi(\bar c,\bar x)\ \Longleftrightarrow\ \forall\bar c.(\bar c\#\bar x\Rightarrow\phi(\bar c,\bar x))\ \Longleftrightarrow\ \exists\bar c.(\bar c\#\bar x\land\phi(\bar c,\bar x)).
\]
\begin{proof}
By induction on the length of $\bar c$. The base case is (4).
For the inductive step, write $\bar c'=\bar c,c_{n+1}$ and define $\psi(\bar c,\bar x):=\newq c_{n+1}.\phi(\bar c,c_{n+1},\bar x)$.
Then $\newq\bar c'.\phi\equiv\newq\bar c.\psi$ and by the induction hypothesis and (4) this is equivalent to $\forall\bar c\ \forall c_{n+1}.(\bar c\#\bar x\land c_{n+1}\#\bar x,\bar c\Rightarrow\phi)$.
By Notation A.2, this is exactly $\forall\bar c'.(\bar c'\#\bar x\Rightarrow\phi)$, as required. The existential case is analogous.
\end{proof}

\paragraph{Theorem 2.3.5 (Generalised Fix-Point).} Let $x\in\mathcal X$ and let $\pi$ be a permutation.
\begin{enumerate}
\item If $\pi$ is a $\bar c$-permutation and $A=\operatorname{mv}(\pi)\setminus\{\bar c\}$, then $\newq\bar c.\pi\cdot x=x$ iff $A\#x$.
\item $\newq\bar c.\pi\cdot x=x$ iff $(\newq\bar c.\pi_{\bar c}\cdot x=x)\land(\pi_{\neg\bar c}\cdot x=x)$.
Equivalently, if $A=\operatorname{mv}(\pi_{\bar c})\setminus\{\bar c\}$, then $\newq\bar c.\pi\cdot x=x$ iff $A\#x\land(\pi_{\neg\bar c}\cdot x=x)$.
\end{enumerate}

\begin{proof}
\begin{enumerate}
\item $(\Rightarrow)$ Assume $\newq\bar c.\pi\cdot x=x$.
By Lemma A.3, for some (equivalently, any) choice of $\bar c$ such that $\bar c\#x$, we have $\pi\cdot x=x$.
Let $\eta$ be a cycle of $\pi$. Since $\pi$ is a $\bar c$-permutation, $\eta$ contains some atom $a\in\bar c$.
As $a\#x$, Lemma 1 yields $\operatorname{mv}(\eta)\#x$. Hence $\operatorname{mv}(\pi)\#x$, and therefore $A\#x$.
$(\Leftarrow)$ Assume $A\#x$. Let $\bar c$ be such that $\bar c\#x$.
Since $\pi$ is a $\bar c$-permutation, every cycle of $\pi$ contains an atom from $\bar c$ and all remaining atoms lie in $A$.
Thus every cycle is disjoint from $\operatorname{supp}(x)$, so $\pi\cdot x=x$.
By Lemma A.3, it follows that $\newq\bar c.\pi\cdot x=x$.
\item If $\pi_{\bar c}=\mathsf{id}$ or $\pi_{\neg\bar c}=\mathsf{id}$, the result is immediate. Assume both are non-identity.
We first show that $\newq\bar c.\pi\cdot x=x$ iff $\newq\bar c.(\pi_{\bar c}\cdot x=x\land\pi_{\neg\bar c}\cdot x=x)$.
By Lemma A.3, both sides are equivalent to universal statements over $\bar c$ with the freshness condition $\bar c\#x,\operatorname{mv}(\pi)\setminus\{\bar c\}$.
Hence it suffices to show that, for such $\bar c$, one has $\pi\cdot x=x$ iff both $\pi_{\bar c}\cdot x=x$ and $\pi_{\neg\bar c}\cdot x=x$.
Assume $\pi\cdot x=x$. Let $\eta$ be a cycle of $\pi_{\bar c}$. Since $\pi_{\bar c}$ is a $\bar c$-permutation, $\eta$ contains some atom $a\in\bar c$.
As $a\#x$, Lemma 1 yields $\operatorname{mv}(\eta)\#x$. Since this holds for every cycle $\eta$ of $\pi_{\bar c}$, and the cycles are disjoint, we obtain $\operatorname{mv}(\pi_{\bar c})\#x$, and therefore $\pi_{\bar c}\cdot x=x$ (equivalently, $\pi_{\bar c}^{-1}\cdot x=x$).
From $\pi=\pi_{\bar c}\circ\pi_{\neg\bar c}$ and $\pi\cdot x=x$, we have $\pi_{\bar c}\cdot(\pi_{\neg\bar c}\cdot x)=x$.
Applying $\pi_{\bar c}^{-1}$ to both sides and using that $\pi_{\bar c}^{-1}\cdot x=x$, it follows that $\pi_{\neg\bar c}\cdot x=x$.
Conversely, if both $\pi_{\bar c}$ and $\pi_{\neg\bar c}$ fix $x$, then $\pi\cdot x=(\pi_{\bar c}\circ\pi_{\neg\bar c})\cdot x=\pi_{\bar c}\cdot(\pi_{\neg\bar c}\cdot x)=x$.
From the fact that $\mathsf N$ distributes over conjunction, we obtain $\newq\bar c.(\pi_{\bar c}\cdot x=x\land\pi_{\neg\bar c}\cdot x=x)$ is equivalent to $\newq\bar c.(\pi_{\bar c}\cdot x=x)\land\newq\bar c.(\pi_{\neg\bar c}\cdot x=x)$.
Since $\pi_{\neg\bar c}$ does not involve atoms from $\bar c$, the formula $\pi_{\neg\bar c}\cdot x=x$ is independent of $\bar c$, and hence, $\newq\bar c.(\pi_{\neg\bar c}\cdot x=x)$ is equivalent to $\pi_{\neg\bar c}\cdot x=x$.
Therefore, $\newq\bar c.\pi\cdot x=x$ iff $\newq\bar c.(\pi_{\bar c}\cdot x=x)\land(\pi_{\neg\bar c}\cdot x=x)$, and the final equivalence follows from item (1), applied to $\pi_{\bar c}$ and $B$.
\end{enumerate}
\end{proof}

\begin{thebibliography}{99}
\bibitem[1]{source1} Mauricio Ayala-Rincón, Maribel Fernández, and Daniele Nantes-Sobrinho. On Nominal Syntax and Permutation Fixed Points. Log. Methods Comput. Sci., 16(1):19:1-19:36, 2020. URL: https://doi.org/10.23638/LMCS-16(1:19)2020 .
\bibitem[2]{source2} Franz Baader and Tobias Nipkow. Term Rewriting and All That. Cambridge University Press, 1998. URL: https://doi.org/10.1017/CBO9781139172752 .
\bibitem[3]{source3} Mikolaj Bojanczyk, Bartek Klin, and Slawomir Lasota. Automata theory in nominal sets. Log. Methods Comput. Sci., 10(3), 2014. URL: https://doi.org/10.2168/LMCS-10(3:4)2014 .
\bibitem[4]{source4} Stanley Burris and Hanamantagouda P. Sankappanavar. A Course in Universal Algebra, volume 78 of Graduate Texts in Mathematics. Springer, 1981. URL: https://link.springer.com/book/9781461381327 .
\bibitem[5]{source5} Ali K. Caires-Santos, Maribel Fernández, and Daniele Nantes-Sobrinho. Equational Reasoning Modulo Commutativity in Languages with Binders. In Automated Deduction – CADE 30: 30th International Conference on Automated Deduction, Stuttgart, Germany, July 28-31, 2025, Proceedings, pages 632-652, Berlin, Heidelberg, 2025. Springer-Verlag. URL: https://doi.org/10.1007/978-3-031-99984-0\_33 .
\bibitem[6]{source6} Ali K. Caires-Santos, Maribel Fernández, and Daniele Nantes-Sobrinho. Strong Nominal Semantics for Fixed-Point Constraints. Electronic Notes in Theoretical Informatics and Computer Science, Proceedings of MFPS XL., 4, December 2024. URL: https://doi.org/10.46298/entics.14777 .
\bibitem[7]{source7} Arthur Charguéraud. The locally nameless representation. J. Autom. Reason., 49(3):363-408, 2012. URL: https://doi.org/10.1007/s10817-011-9225-2 .
\bibitem[8]{source8} Manuel Clavel, Francisco Durán, Steven Eker, Patrick Lincoln, Narciso Martí-Oliet, José Meseguer, and Carolyn L. Talcott, editors. All About Maude - A High-Performance Logical Framework, volume 4350 of LNCS. Springer, 2007. URL: https://doi.org/10.1007/978-3-540-71999-1 .
\bibitem[9]{source9} N.G de Bruijn. Lambda calculus notation with nameless dummies, a tool for automatic formula manipulation, with application to the church-rosser theorem. Indagationes Mathematicae (Proceedings), 75(5):381-392, 1972. URL: https://doi.org/10.1016/1385-7258(72)90034-0 .
\bibitem[10]{source10} Francisco Durán, Santiago Escobar, Jose Meseguer, and Julia Sapiña. Nuitp: An inductive theorem prover for equational program verification. In Proceedings of the 26th International Symposium on Principles and Practice of Declarative Programming, PPDP '24, New York, NY, USA, 2024. Association for Computing Machinery. URL: https://doi.org/10.1145/3678232.3678236 .
\bibitem[11]{source11} Amy P. Felty, Alberto Momigliano, and Brigitte Pientka. The next 700 challenge problems for reasoning with higher-order abstract syntax representations - part 2 - A survey. J. Autom. Reason., 55(4):307-372, 2015. URL: https://doi.org/10.1007/s10817-015-9327-3 .
\bibitem[12]{source12} Maribel Fernández and Murdoch James Gabbay. Nominal Rewriting. Inf. Comput., 205(6):917-965, 2007. URL: https://doi.org/10.1016/j.ic.2006.12.002 .
\bibitem[13]{source13} Murdoch James Gabbay. Nominal Algebra and the HSP Theorem. J. Log. Comput., 19(2):341-367, 2009. URL: https://doi.org/10.1093/logcom/exn055 .
\bibitem[14]{source14} Murdoch James Gabbay. Foundations of Nominal Techniques: Logic and Semantics of Variables in Abstract Syntax. Bull. Symb. Log., 17(2):161-229, 2011. URL: https://doi.org/10.2178/bsl/1305810911 .
\bibitem[15]{source15} Murdoch James Gabbay and Aad Mathijssen. Nominal (Universal) Algebra: Equational Logic with Names and Binding. J. Log. Comput., 19(6):1455-1508, 2009. URL: https://doi.org/10.1093/logcom/exp033 .
\bibitem[16]{source16} Murdoch James Gabbay and Andrew M. Pitts. A New Approach to Abstract Syntax with Variable Binding. Formal Aspects Comput., 13(3-5):341-363, 2002. URL: https://doi.org/10.1007/s001650200016 .
\bibitem[17]{source17} Alexander Kurz and Daniela Petrişan. On universal algebra over nominal sets. Mathematical Structures in Computer Science, 20(2):285-318, 2010. URL: https://doi.org/10.1017/S0960129509990399 .
\bibitem[18]{source18} José Meseguer. Inductive reasoning with equality predicates, contextual rewriting and variant-based simplification. J. Log. Algebraic Methods Program., 144:101036, 2025. URL: https://doi.org/10.1016/j.jlamp.2025.101036 .
\bibitem[19]{source19} Robin Milner, Joachim Parrow, and David Walker. A Calculus of Mobile Processes, I. Information and Computation, 100(1):1-40, 1992. URL: https://doi.org/10.1016/0890-5401(92)90008-4 .
\bibitem[20]{source20} Peter W. O'Hearn, John C. Reynolds, and Hongseok Yang. Local reasoning about programs that alter data structures. In Laurent Fribourg, editor, Computer Science Logic (CSL 2001), volume 2142 of Lecture Notes in Computer Science, pages 1-19, Berlin, Heidelberg, 2001. Springer. URL: https://doi.org/10.1007/3-540-44802-0\_1 .
\bibitem[21]{source21} F. Pfenning and C. Elliott. Higher-Order Abstract Syntax. SIGPLAN Not., 23(7):199-208, 1988. URL: https://doi.org/10.1145/960116.54010 .
\bibitem[22]{source22} Andrew M. Pitts. Nominal Sets: Names and Symmetry in Computer Science. Cambridge Tracts in Theoretical Computer Science. Cambridge University Press, 2013. URL: https://doi.org/10.1017/CBO9781139084673 .
\bibitem[23]{source23} Andrei Popescu. Nominal recursors as epi-recursors. Proc. ACM Program. Lang., 8(POPL):425-456, 2024. URL: https://doi.org/10.1145/3632857 .
\bibitem[24]{source24} John C. Reynolds. Separation logic: A logic for shared mutable data structures. In Proceedings of the 17th Annual IEEE Symposium on Logic in Computer Science (LICS 2002), pages 55-74, Washington, DC, USA, 2002. IEEE Computer Society. URL: https://doi.org/10.1109/LICS.2002.1029817 .
\bibitem[25]{source25} Joseph J. Rotman. An Introduction to the Theory of Groups, volume 148 of Graduate Texts in Mathematics. Springer, New York, 4 edition, 1995. URL: https://link.springer.com/book/10.1007/978-1-4612-4176-8 .
\bibitem[26]{source26} Christian Urban, Andrew M. Pitts, and Murdoch James Gabbay. Nominal Unification. Theor. Comput. Sci., 323(1-3):473-497, 2004. URL: https://doi.org/10.1016/j.tcs.2004.06.016 .
\bibitem[27]{source27} Jan van Brügge, James McKinna, Andrei Popescu, and Dmitriy Traytel. Barendregt convenes with Knaster and Tarski: Strong rule induction for syntax with bindings. Proc. ACM Program. Lang., 9(POPL):1687-1718, 2025. URL: https://doi.org/10.1145/3704893 .
\end{thebibliography}
\end{document}
